% !TEX root = ../main.tex
\chapter{Reproducibility and Determinism}
\label{ch:reproducibility}

The final chapter treats reproducibility: the framework
property that ensures the same code, run on the same data,
produces the same numbers. Reproducibility is a precondition
for debugging, scientific publication, and CI; it is also one
of the hardest properties to deliver in a parallel framework
because most parallel primitives are non-associative in
floating point. This chapter documents \Lucid's
reproducibility guarantees (what is deterministic, what is
not, by-design), the global determinism mode that trades
performance for bit-exact runs, the random-number-generation
infrastructure, the dataset-side reproducibility guarantees,
and the cross-version stability story. The
\cite{pineau_reproducibility_2021} report is the field's
reference for the broader problem; this chapter is \Lucid's
specific implementation.

\section{The Reproducibility Spectrum}
\label{sec:reproducibility:spectrum}

``Reproducible'' is not binary. Useful levels:

\begin{enumerate}
  \item \textbf{Bit-exact}: two runs produce numerically
    identical tensors. The strongest guarantee; expensive.
  \item \textbf{Trajectory-equivalent}: two runs follow the
    same training trajectory (loss within $10^{-6}$ at every
    step). Useful for debugging.
  \item \textbf{Outcome-equivalent}: two runs reach the same
    final accuracy ($\pm$\,0.1\,\%). Useful for paper
    reproduction.
  \item \textbf{Statistical}: two runs produce results from
    the same distribution. The weakest guarantee; useful for
    research.
\end{enumerate}

\Lucid\ targets bit-exact for fixed seed + fixed hardware +
deterministic mode. Trajectory-equivalent is the default
(no special mode); outcome-equivalent is automatic across
similar hardware.

\section{Sources of Non-Determinism}
\label{sec:reproducibility:sources}

Non-determinism arises from:

\begin{itemize}
  \item \textbf{Floating-point non-associativity}: parallel
    reductions sum in different orders, producing different
    results in the last digits.
  \item \textbf{GPU kernel scheduling}: the order in which a
    matmul's tiles are assigned to GPU cores can vary; combined
    with non-associativity, this gives different totals.
  \item \textbf{Memory allocator state}: an allocator may
    return different memory addresses on different runs;
    if the addresses are read as part of a hash (unlikely but
    possible), the result differs.
  \item \textbf{Multi-threaded reductions}: CPU code with
    OpenMP-style reductions sums in different orders.
  \item \textbf{Worker-process ordering}: data loaders with
    multiple workers may deliver batches in nondeterministic
    order if the queueing is unfair.
  \item \textbf{System time-based RNG}: any code that calls
    \code{time.time()} as a seed source is non-reproducible.
\end{itemize}

\Lucid's determinism mode addresses each.

\section{The Determinism Mode}
\label{sec:reproducibility:det_mode}

The user opts in via:

\begin{lstlisting}[style=lucid,language=LucidPython]
lucid.use_deterministic_algorithms(True)
\end{lstlisting}

When enabled:

\begin{itemize}
  \item Reductions use a fixed serial order (slower, bit-exact).
  \item GPU kernels prefer deterministic variants where
    available; raise on non-deterministic ones.
  \item The allocator uses a fixed-seed strategy (allocations
    in the same logical order always return the same physical
    addresses).
  \item Multi-threaded CPU code uses serial paths or
    fixed-order parallel paths.
\end{itemize}

The performance cost is typically 1.5--3$\times$ on the
affected workloads.

\subsection{What raises in determinism mode}
\label{ssec:reproducibility:det_raises}

If a non-deterministic op is called without a deterministic
alternative, \Lucid\ raises:

\begin{lstlisting}[style=lucid,language=LucidPython]
DeterministicError: scatter_add with float inputs uses a
  non-deterministic GPU kernel (parallel sums). No
  deterministic alternative is available in 3.5.
  Workaround: move to CPU stream (slower but deterministic).
\end{lstlisting}

The user can either accept the slowdown of the CPU path or
opt out of strict determinism for that subgraph.

\section{Random Number Generation}
\label{sec:reproducibility:rng}

The RNG is the most user-visible source of non-determinism.
\Lucid's RNG is a deterministic counter-based generator
(Philox \cite{philox_salmon_2011}) that produces the same
sequence given the same seed:

\begin{lstlisting}[style=lucid,language=LucidPython]
lucid.seed(42)
a = lucid.randn(100)

lucid.seed(42)
b = lucid.randn(100)

assert (a == b).all().item()  # True, always
\end{lstlisting}

\subsection{Why Philox, not Mersenne Twister}
\label{ssec:reproducibility:why_philox}

The Mersenne Twister \cite{matsumoto_mt_1998} (the standard
Python RNG) is sequential: parallel generation requires
splitting into substreams, which is awkward. Philox is
counter-based: any subset of the sequence can be generated
in parallel without communication. For a tensor-generating
RNG, this matters.

\subsection{Per-stream state}
\label{ssec:reproducibility:per_stream_state}

\Lucid\ maintains separate RNG state for the CPU and GPU
streams. \code{lucid.seed(42)} seeds both.

\begin{lstlisting}[style=lucid,language=LucidPython]
state = lucid.random.get_rng_state()
# (cpu_state, gpu_state)

# ... some work ...

lucid.random.set_rng_state(state)
# Now both streams are restored to the saved point
\end{lstlisting}

\subsection{Generator objects}
\label{ssec:reproducibility:generators}

For finer-grained control, users can create per-op generators:

\begin{lstlisting}[style=lucid,language=LucidPython]
gen = lucid.Generator()
gen.manual_seed(42)
a = lucid.randn(100, generator=gen)
b = lucid.randn(100, generator=gen)
# a and b are different (gen has stepped), but reproducible
\end{lstlisting}

\section{Worker-Process Determinism}
\label{sec:reproducibility:worker_det}

For multi-process data loaders, each worker gets a deterministic
per-worker seed:

\begin{lstlisting}[style=lucid,language=LucidPython]
def worker_init_fn(worker_id):
    lucid.seed(42 + worker_id)

loader = DataLoader(
    dataset,
    num_workers=4,
    worker_init_fn=worker_init_fn,
    generator=lucid.Generator().manual_seed(42),
)
\end{lstlisting}

The \code{generator} parameter controls the per-epoch shuffle;
\code{worker\_init\_fn} controls per-worker augmentation
randomness. Together, two runs with the same seed produce
identical batch sequences.

\subsection{Worker output ordering}
\label{ssec:reproducibility:worker_ordering}

The data loader can deliver batches in two orders:
\begin{itemize}
  \item \textbf{As they finish} (default): better throughput,
    non-deterministic order (different runs may see batches in
    different orders).
  \item \textbf{In submitted order} (\code{ordered=True}): the
    loader buffers out-of-order batches until the next-in-line
    is ready. Deterministic, slightly slower.
\end{itemize}

For strict reproducibility, \code{ordered=True}.

\section{Cross-Hardware Determinism}
\label{sec:reproducibility:cross_hardware}

Bit-exact determinism across different hardware is generally
infeasible: M1 Pro and M3 Max have different GPU core counts
and may make different kernel-scheduling decisions even in
deterministic mode. \Lucid's stance:

\begin{itemize}
  \item Same hardware + same seed + determinism mode:
    bit-exact, guaranteed.
  \item Different hardware (M1 vs M3): trajectory-equivalent
    (loss within $10^{-5}$).
  \item Different OS version: trajectory-equivalent if the
    same MPSGraph feature set is available.
\end{itemize}

For paper-reproduction purposes, this is sufficient; the
training trajectory of a model is stable across the M-series
family.

\section{Cross-Version Reproducibility}
\label{sec:reproducibility:cross_version}

A model trained on \Lucid\ 3.5.0 should reproduce on \Lucid\
3.5.x. \Lucid's policy:

\begin{itemize}
  \item Within a minor version (3.5.0, 3.5.1, 3.5.2):
    bit-exact in determinism mode for the same model and data.
  \item Across minor versions (3.5 to 3.6): trajectory-
    equivalent; the model converges to within statistical
    noise.
  \item Across major versions (3.x to 4.x): not guaranteed;
    refactors may change numerical paths.
\end{itemize}

The compatibility tests (\code{lucid/test/reproducibility/})
verify the first two guarantees.

\section{Checkpoint Reproducibility}
\label{sec:reproducibility:checkpoint}

A checkpoint saved during training contains:
\begin{itemize}
  \item Model state\_dict.
  \item Optimiser state.
  \item Scheduler state.
  \item RNG state (CPU + GPU).
  \item Global step count.
  \item Random seeds (for resume).
\end{itemize}

The standard save:

\begin{lstlisting}[style=lucid,language=LucidPython]
lucid.save({
    "model":      model.state_dict(),
    "optimizer":  optimizer.state_dict(),
    "scheduler":  scheduler.state_dict(),
    "rng":        lucid.random.get_rng_state(),
    "step":       global_step,
    "seed":       42,  # The seed used at training start
}, "checkpoint.lpt")
\end{lstlisting}

Resume:

\begin{lstlisting}[style=lucid,language=LucidPython]
ckpt = lucid.load("checkpoint.lpt")
lucid.seed(ckpt["seed"])
model.load_state_dict(ckpt["model"])
optimizer.load_state_dict(ckpt["optimizer"])
scheduler.load_state_dict(ckpt["scheduler"])
lucid.random.set_rng_state(ckpt["rng"])
global_step = ckpt["step"]
\end{lstlisting}

After this, training continues bit-exactly from where it
stopped (assuming the same hardware + same data + deterministic
mode).

\section{Environment Capture}
\label{sec:reproducibility:env_capture}

For paper reproduction, the user should capture the full
environment:

\begin{lstlisting}[style=lucid,language=LucidPython]
env = lucid.utils.repro.capture_environment()
# {
#   "lucid_version": "3.5.0",
#   "mlx_version": "0.31.2",
#   "macos_version": "26.0.1",
#   "hardware": {"chip": "Apple M3 Max", "ram_gb": 64, ...},
#   "python_version": "3.14.0",
#   "git_commit": "<repo SHA at time of capture>",
#   "deterministic_mode": True,
#   "rng_seed": 42,
# }
lucid.save(env, "experiment_env.json")
\end{lstlisting}

The captured environment file should be saved alongside the
trained model and published with the paper.

\section{Reproducibility Checks in CI}
\label{sec:reproducibility:ci_check}

\Lucid's CI includes a small reproducibility test suite:

\begin{lstlisting}[style=lucid,language=LucidPython]
def test_resnet50_bit_exact():
    """Two runs with same seed should produce identical loss."""
    lucid.use_deterministic_algorithms(True)
    lucid.seed(42)
    losses1 = [train_step(...) for _ in range(20)]

    lucid.seed(42)
    losses2 = [train_step(...) for _ in range(20)]

    for l1, l2 in zip(losses1, losses2):
        assert l1 == l2, f"Determinism failure: {l1} vs {l2}"
\end{lstlisting}

The test runs on every release-candidate build. Any commit
that breaks bit-exact determinism fails the test (the author
must either restore determinism or document why the break is
acceptable).

\section{The Limits of Determinism}
\label{sec:reproducibility:limits}

Despite the best efforts, some operations are fundamentally
non-deterministic (or only deterministic at unacceptable cost):

\begin{itemize}
  \item \textbf{Float reductions across many cores}:
    bit-exact only with serial reduction.
  \item \textbf{Atomic float updates}: e.g., \code{scatter\_add}
    on overlapping indices.
  \item \textbf{Some MPSGraph fused kernels}: Apple does not
    guarantee bit-exact reductions in all kernels.
\end{itemize}

For these, \Lucid\ documents the non-determinism explicitly in
the op's docstring; users opting into determinism mode get a
clear error rather than silent variance.

\section{Reproducibility vs Performance}
\label{sec:reproducibility:tradeoff}

The cost of full determinism:

\begin{table}[t]
\centering
\small
\begin{tabular}{lrrr}
\toprule
Workload & Non-det (ms) & Det (ms) & Slowdown \\
\midrule
ResNet-50 train       & 64  & 95  & 1.48$\times$ \\
GPT-2 small train     & 105 & 175 & 1.67$\times$ \\
BERT-base train       & 88  & 130 & 1.48$\times$ \\
\bottomrule
\end{tabular}
\caption[Determinism cost.]{Iteration time with vs without
strict determinism. The slowdown comes from serial reductions
in attention's softmax and the optimiser's parameter updates
(both have multi-core reductions in the fast path).}
\label{tab:reproducibility:cost}
\end{table}

For day-to-day training, non-deterministic mode is the default;
determinism is opted-into for debugging or publication runs.

\section{The Reproducibility Pledge}
\label{sec:reproducibility:pledge}

\Lucid's commitment to users:

\begin{enumerate}
  \item Any \Lucid\ program with a fixed seed + deterministic
    mode + fixed hardware produces bit-exact output.
  \item Any minor-version release maintains this property for
    pre-existing programs.
  \item Sources of inherent non-determinism are documented and
    raise in determinism mode if no deterministic alternative
    exists.
  \item The reproducibility tests are CI-enforced; PRs that
    break them are rejected or must justify the change.
\end{enumerate}

\section{Summary and Forward References}
\label{sec:reproducibility:summary}

\Lucid's reproducibility infrastructure: opt-in determinism
mode, Philox-based RNG with per-stream state, worker-deterministic
data loading, checkpoint format that captures all state needed
for bit-exact resume, environment-capture utility for paper
publication, CI checks that enforce bit-exact for the
deterministic mode. The performance cost of determinism is
1.5$\times$; the user chooses when to pay it.

This chapter closes \partref{part:production} and the report.
The report has traversed Apple silicon hardware, the framework's
layered architecture, every op family, every layer, the model
zoo, the performance work, and the production-engineering
support stack. A user finishing the report has a complete
picture of \Lucid: what it is, why it is, how to use it, and
how to extend it.

\begin{lucidnote}
For the user: for paper-reproducible runs, enable
\code{lucid.use\_deterministic\_algorithms(True)} and
fix every random seed (numpy, Python, \Lucid). For day-to-day
training, leave determinism off; the trajectory is stable
enough that final accuracy is reproducible to $\pm$0.1\%.
\end{lucidnote}
